\documentclass[11pt,a4paper]{article}
% Preámbulo compartido para los modelos del segundo parcial de ACN.
% Replica el formato del parcial real 2P-2S-2025.
\usepackage[utf8]{inputenc}
\usepackage[T1]{fontenc}
\usepackage[spanish]{babel}
\usepackage{amsmath,amssymb}
\usepackage[margin=2.6cm]{geometry}
\usepackage{enumitem}
\usepackage{array}

\pagestyle{plain}
\setlength{\parindent}{0pt}
\setlength{\parskip}{5pt}

% Listas de incisos: a), b), c), ...
\setlist[enumerate,1]{label=\alph*),leftmargin=2.2em,itemsep=7pt,topsep=7pt}
% Listas numeradas 1., 2., 3., ...
\newlist{puntos}{enumerate}{1}
\setlist[puntos]{label=\arabic*.,leftmargin=2.2em,itemsep=7pt,topsep=7pt}
% Listas de condiciones del enunciado
\setlist[itemize,1]{label=$\bullet$,leftmargin=2em,itemsep=5pt,topsep=6pt}

% Encabezado del parcial
\newcommand{\encabezado}[1]{%
  \begin{center}
    \textbf{\large Aplicaciones Computacionales en Negocios}\\[3pt]
    \textbf{\large Segundo semestre de 2025 --- Segundo parcial}\\[3pt]
    \textbf{\normalsize Modelo de práctica #1}
  \end{center}
  \vspace{4pt}
  \hrule height 0.8pt
  \vspace{14pt}
  Nombre:\par\vspace{12pt}
  Legajo:\par\vspace{12pt}
  DNI:\par\vspace{12pt}
  ¿Te recibís con este parcial?:
  \vspace{10pt}
  \hrule height 0.8pt
  \vspace{16pt}
}

% Título de ejercicio
\newcommand{\ejercicio}[1]{\par\medskip\noindent\textbf{Ejercicio #1.}\ }

% Atajos de notación usados por la cátedra
\newcommand{\R}{\mathbb{R}}
\newcommand{\Z}{\mathbb{Z}}
\newcommand{\dem}{\mathrm{dem}}
\newcommand{\capa}{\mathrm{cap}}
\newcommand{\costo}{\mathrm{costo}}


\begin{document}
\encabezado{4}

%======================================================================
\ejercicio{1}
Consideremos el siguiente modelo de \emph{lot sizing} con un único producto, en
el que se permite atrasar el cumplimiento de la demanda (\emph{backlogging}).
Tenemos un conjunto $T = \{1,\dots,n\}$ de períodos; para cada $t \in T$ tenemos
la demanda $d_t \in \R_+$, el costo unitario de producción $c_t \in \R_+$, el
costo unitario de almacenamiento $a_t \in \R_+$ y el costo $b_t \in \R_+$ por
cada unidad de demanda atrasada al final del período $t$. La variable
$x_t \in \R_+$ representa la cantidad fabricada en el período $t$, la variable
$s_t \in \R_+$ el stock al final del período $t$, y la variable $r_t \in \R_+$ la
demanda acumulada que todavía no fue atendida al final del período $t$:

\begin{align}
\min \quad & \sum_{t \in T} \big( c_t\, x_t + a_t\, s_t + b_t\, r_t \big) \\
s_t - r_t &= s_{t-1} - r_{t-1} + x_t - d_t && \forall\, t \in T \\
x_t,\ s_t,\ r_t &\ge 0 && \forall\, t \in T
\end{align}

con $s_0 = r_0 = 0$.

\begin{enumerate}
  \item Explicar por qué, si $a_t + b_t > 0$ para todo $t \in T$, en ninguna
  solución óptima puede suceder que $s_t > 0$ y $r_t > 0$ al mismo tiempo.
  ¿Qué sucede con las soluciones óptimas si $b_t = 0$ para todo $t \in T$?
  ¿Qué restricción hay que agregar para garantizar que toda la demanda quede
  atendida al final del horizonte?

  \item Supongamos que encender la línea en el período $t$ tiene un costo fijo
  $f_t \in \R_+$, y que no hay una capacidad de producción explícita. ¿Cómo se
  modifica el modelo? Proponer un valor concreto para la constante $M$ de la
  restricción que liga la producción con el encendido de la línea. ¿Por qué el
  valor que se suele usar en el modelo \emph{sin} backlogging ya no es
  correcto? ¿Por qué conviene que $M$ sea lo más chico posible?

  \item Agregar al modelo del punto anterior una restricción para que ninguna
  unidad de demanda se atienda con más de un período de atraso (suponiendo que
  la demanda atrasada se atiende en orden de llegada). Agregar además una
  restricción para que haya demanda atrasada en a lo sumo tres períodos de todo
  el horizonte.

  \item Supongamos ahora que, además del costo fijo $f_t$, se paga un costo de
  arranque $g \in \R_+$ cada vez que la línea se enciende en un período $t$
  habiendo estado apagada en el período $t-1$ (la línea está apagada en el
  período 0). ¿Cómo se agrega este costo al modelo? ¿Hace falta que las
  variables nuevas sean binarias? Justificar.
\end{enumerate}

\newpage
%======================================================================
\ejercicio{2}
Estamos a cargo de la dotación de un \emph{call center} que atiende las 24
horas. Sea $H = \{0, 1, \dots, 23\}$ el conjunto de las franjas horarias de un
día (la franja $h$ empieza a la hora $h$ y termina a la hora $h+1$), y para cada
$h \in H$ conocemos la cantidad mínima $r_h \in \Z_+$ de operadores que deben
estar trabajando en esa franja. La planificación se repite todos los días, de
modo que un turno que empieza a las 22 continúa en las primeras horas del día
siguiente. Hay dos tipos de operadores:

\begin{itemize}
  \item Los operadores \emph{full-time} trabajan 8 horas consecutivas y cobran
  $c^F$ pesos por turno.
  \item Los operadores \emph{part-time} trabajan 4 horas consecutivas y cobran
  $c^P$ pesos por turno.
  \item Cada turno (de cualquiera de los dos tipos) puede empezar al comienzo de
  cualquier franja $h \in H$.
  \item Los turnos que incluyen alguna franja entre las 0 y las 6 (es decir, las
  franjas $0, \dots, 5$) tienen un adicional nocturno de $\nu$ pesos por
  operador, sin importar el tipo de turno.
\end{itemize}

\begin{puntos}
  \item Formular un modelo de programación lineal entera para determinar
  cuántos operadores de cada tipo comienzan su turno en cada franja, de modo
  tal de cubrir los requerimientos a costo mínimo. \emph{Pista:} usar variables
  enteras $F_h$ y $P_h$ que representen la cantidad de operadores full-time y
  part-time que comienzan su turno en la franja $h$. Escribir explícitamente la
  expresión de la cantidad de operadores activos en la franja $h$, teniendo en
  cuenta el carácter cíclico del día.

  \item Agregar al modelo las siguientes condiciones: (i) la cantidad total de
  operadores part-time no puede superar la mitad de la cantidad total de
  operadores full-time; (ii) ningún operador part-time puede comenzar su turno
  entre las 0 y las 6.

  \item Para simplificar la gestión, la empresa quiere usar a lo sumo cinco
  horarios distintos de inicio para los turnos full-time y, si en un horario
  comienzan operadores full-time, deben comenzar al menos tres. Agregar estas
  condiciones al modelo. Proponer un valor concreto de la constante $M$ que se
  utilice, en función de los datos.

  \item Supongamos que resolvemos la relajación lineal del modelo del primer
  punto. ¿Qué relación hay entre su valor óptimo y el valor óptimo del modelo
  entero? Si redondeamos hacia arriba todas las variables de la solución de la
  relajación, ¿obtenemos una solución factible? ¿Y una solución óptima? ¿Siguen
  valiendo estas respuestas para el modelo del punto anterior? Justificar.
\end{puntos}

\newpage
%======================================================================
\ejercicio{3}
Debemos decidir dónde ubicar las bases de ambulancias de una ciudad. Tenemos un
conjunto $L$ de bases candidatas y un conjunto $C$ de zonas. Para cada zona
$i \in C$ conocemos su población $p_i \in \Z_+$ y, para cada base $j \in L$, el
tiempo de viaje $t_{ij} \in \R_+$ entre la base $j$ y la zona $i$. Abrir la base
$j$ tiene un costo $f_j \in \R_+$. Decimos que la base $j$ \emph{cubre} a la zona
$i$ si $t_{ij} \le U$, donde $U$ es un dato (por ejemplo, 10 minutos), y
llamamos $N_i = \{ j \in L : t_{ij} \le U \}$ al conjunto de bases que cubren a
la zona $i$.

\begin{puntos}
  \item Formular un modelo de programación lineal entera para determinar qué
  bases abrir, de modo tal que todas las zonas queden cubiertas y se minimice el
  costo total de apertura. ¿Puede este modelo ser infactible? ¿Cómo se detecta
  esta situación a partir de los datos, sin resolver el modelo?

  \item Supongamos ahora que tenemos un presupuesto $B$ y que, en lugar de
  exigir que todas las zonas queden cubiertas, queremos maximizar la población
  cubierta. Formular el modelo correspondiente. \emph{Pista:} agregar una
  variable $z_i$ por cada zona $i \in C$ que indique si la zona queda cubierta.
  Explicar por qué alcanza con una sola desigualdad por zona para ligar $z_i$
  con las bases abiertas.

  \item Disponemos de $K$ ambulancias en total, y cada base abierta $j$ puede
  alojar entre 1 y $k_j$ ambulancias. Las zonas del conjunto $C' \subseteq C$
  (las de mayor población) deben quedar cubiertas por al menos dos ambulancias,
  aunque estén en la misma base. Modificar el modelo del punto anterior para
  incorporar estas condiciones.

  \item Consideremos una instancia del modelo del primer punto con tres zonas
  $\{1,2,3\}$ y tres bases $\{A,B,C\}$, todas con costo 1, donde
  $N_1 = \{A, B\}$, $N_2 = \{B, C\}$ y $N_3 = \{A, C\}$. Mostrar que la
  relajación lineal tiene una solución con valor $3/2$, mientras que el óptimo
  entero vale 2. Proponer una desigualdad válida que elimine esa solución
  fraccionaria, y explicar cómo se usaría este tipo de desigualdades dentro de un
  algoritmo de \emph{branch and cut}.
\end{puntos}

\end{document}
